\subsection{Patch Generations}
\label{appx:analyses:patch_generations}

In this section, we present some statistics and analysis around the edits generated by SWE-agent.
At the end of a task episode, the edits made by SWE-agent are aggregated and saved as a single \texttt{.patch} file, the canonical representation for code changes of a pull request on GitHub.
From these patch representations, we can quantitatively characterize an agent's generations and see how they compare to the original solutions written by human codebase maintainers.

Table~\ref{tab.appx_tables.patch_stats_basic} presents a summary of four basic statistics about the model generations.
Lines added and lines removed refer to the total number of lines that were added or deleted in the patch, an indicator of the size of the modification.
The number of hunks and files is more indicative of how many ``regions" of the codebase were modified.
A higher number of hunks and files suggests that there are more distinct, separate places in the codebase where the patch made changes.
For both ``Resolved" and ``All" categories of task instances, models tend to generate ``larger" edits (e.g., more lines added, hunks, and files) than the corresponding gold solution.
Prior RAG baselines in ~\citet{jimenez2024swebench} typically produce smaller edits on average.
The source of this increase for agent-generated solutions can largely be attributed to additional reproduction code.

\begin{table}[ht]
    \caption{
    We show the (median) / (mean) value for several statistics characterizing patch generations.
    We calculate these statistics across two dimensions.
    First, the ``Resolved" / ``All" labels denote whether the patch resolved the issue.
    Second, for the task instances specific to each model, we calculate the same statistics across the gold patches.
    To diminish the effect of outliers, we calculate these statistics based on values falling within within the \num{90}th percentile of the distribution.
    }
    \label{tab.appx_tables.patch_stats_basic}
    \vspace{0.25em}
    \centering
    \begin{tabular}{llcccc}
        \toprule
        Model & Outcome & Lines + & Lines - & Hunks & Files \\
        \midrule
        SWE-agent & Resolved & 3.0 / 5.7 & 1.0 / 1.32 & 1.0 / 1.52 & 1.0 / 1.22 \\
        \smalltab w/ GPT-4 Turbo & Any & 12.0 / 16.58 & 1.0 / 1.35 & 2.0 / 1.83 & 1.0 / 1.53 \\
        \cmidrule(l){3-6}
        Gold    & Resolved & 2.0 / 3.58 & 1.0 / 1.98 & 1.0 / 1.3 & 1.0 / 1.0 \\
                & Any & 7.0 / 11.67 & 2.0 / 4.05 & 2.0 / 2.45 & 1.0 / 1.24 \\
        \midrule
        SWE-agent & Resolved & 3.0 / 5.09 & 1.0 / 1.59 & 1.0 / 1.56 & 1.0 / 1.26 \\
        \smalltab w/ Claude 3 Opus & Any & 11.0 / 15.25 & 1.0 / 1.79 & 2.0 / 2.14 & 2.0 / 1.87 \\
        \cmidrule(l){3-6}
        Gold & Resolved & 3.0 / 3.91 & 1.0 / 1.94 & 1.0 / 1.4 & 1.0 / 1.0 \\
        & Any & 6.0 / 10.68 & 2.0 / 3.61 & 2.0 / 2.22 & 1.0 / 1.13 \\
        \bottomrule
    \end{tabular}
\end{table}

When comparing the ``Resolved" and ``All" categories, we see that successfully resolved edits are relatively smaller than the original distribution.
This trend is consistent with the RAG based solutions; issues that require multiple edits across a codebase remains challenging for agents.